\documentclass[10pt,a4paper]{amsart}
\usepackage{amsmath,amssymb,amsthm,graphicx}
% Retained preamble for the Stacks Project chapter snapshot in this directory.
% The chapter's own preamble.tex is a copy of the shared Stacks Project preamble; the block
% below is sliced MECHANICALLY, line for line, from it. Nothing was added, removed or
% redefined; the wrapper class, the xr-hyper/hyperref wiring and the externaldocument list
% are omitted because the assembler supplies the document class.
%
% lines taken from preamble.tex: 3, 5, 9, 11, 12, 156, 157, 158, ... (51 lines total)
%
% package coverage restored from the chapter's own preamble (2026-10-13): xy, xy-2cell

\usepackage[all]{xy}
\xyoption{2cell}
\UseAllTwocells

\newenvironment{reference}{\comment}{\endcomment}
\newenvironment{slogan}{\comment}{\endcomment}
\newenvironment{history}{\comment}{\endcomment}
\theoremstyle{plain}
\newtheorem{theorem}[subsection]{Theorem}
\newtheorem{proposition}[subsection]{Proposition}
\newtheorem{lemma}[subsection]{Lemma}
\theoremstyle{definition}
\newtheorem{definition}[subsection]{Definition}
\newtheorem{example}[subsection]{Example}
\newtheorem{exercise}[subsection]{Exercise}
\newtheorem{situation}[subsection]{Situation}
\theoremstyle{remark}
\newtheorem{remark}[subsection]{Remark}
\newtheorem{remarks}[subsection]{Remarks}
\numberwithin{equation}{subsection}
\def\lim{\mathop{\mathrm{lim}}\nolimits}
\def\colim{\mathop{\mathrm{colim}}\nolimits}
\def\Spec{\mathop{\mathrm{Spec}}}
\def\Hom{\mathop{\mathrm{Hom}}\nolimits}
\def\Ext{\mathop{\mathrm{Ext}}\nolimits}
\def\SheafHom{\mathop{\mathcal{H}\!\mathit{om}}\nolimits}
\def\SheafExt{\mathop{\mathcal{E}\!\mathit{xt}}\nolimits}
\def\Sch{\mathit{Sch}}
\def\Mor{\mathop{\mathrm{Mor}}\nolimits}
\def\Ob{\mathop{\mathrm{Ob}}\nolimits}
\def\Sh{\mathop{\mathit{Sh}}\nolimits}
\def\NL{\mathop{N\!L}\nolimits}
\def\CH{\mathop{\mathrm{CH}}\nolimits}
\def\proetale{{pro\text{-}\acute{e}tale}}
\def\etale{{\acute{e}tale}}
\def\QCoh{\mathit{QCoh}}
\def\Ker{\mathop{\mathrm{Ker}}}
\def\Im{\mathop{\mathrm{Im}}}
\def\Coker{\mathop{\mathrm{Coker}}}
\def\Coim{\mathop{\mathrm{Coim}}}
\DeclareMathSymbol{\boxtimes}{\mathbin}{AMSa}{"02}
\def\QCohstack{\mathcal{QC}\!\mathit{oh}}
\def\Cohstack{\mathcal{C}\!\mathit{oh}}
\def\Spacesstack{\mathcal{S}\!\mathit{paces}}
\def\Quotfunctor{\mathrm{Quot}}
\def\Hilbfunctor{\mathrm{Hilb}}
\def\Curvesstack{\mathcal{C}\!\mathit{urves}}
\def\Polarizedstack{\mathcal{P}\!\mathit{olarized}}
\def\Complexesstack{\mathcal{C}\!\mathit{omplexes}}
\def\Pic{\mathop{\mathrm{Pic}}\nolimits}
\def\Picardstack{\mathcal{P}\!\mathit{ic}}
\def\Picardfunctor{\mathrm{Pic}}
\def\Deformationcategory{\mathcal{D}\!\mathit{ef}}
\usepackage{cleveref}
\crefname{originalproposition}{original proposition}{original propositions}
\Crefname{originalproposition}{Original proposition}{Original propositions}
\makeatletter
\newlabel{equation-formal-glueing-modules}{{0AEW}{source}{}{sourcepointer-1}{}}
\newlabel{equation-formal-glueing-modules@cref}{{[originalproposition][0][]0AEW}{[1][0][]source}{}{}{}}
\makeatother
\crefname{originalproposition}{original proposition}{original propositions}
\Crefname{originalproposition}{Original proposition}{Original propositions}
\makeatletter
\newlabel{equation-reverse}{{0AEX}{source}{}{sourcepointer-2}{}}
\newlabel{equation-reverse@cref}{{[originalproposition][0][]0AEX}{[1][0][]source}{}{}{}}
\makeatother
\crefname{originalproposition}{original proposition}{original propositions}
\Crefname{originalproposition}{Original proposition}{Original propositions}
\makeatletter
\newlabel{lemma-stalk-formal-glueing}{{0AER}{source}{}{sourcepointer-3}{}}
\newlabel{lemma-stalk-formal-glueing@cref}{{[originalproposition][0][]0AER}{[1][0][]source}{}{}{}}
\makeatother
\crefname{originalproposition}{original proposition}{original propositions}
\Crefname{originalproposition}{Original proposition}{Original propositions}
\makeatletter
\newlabel{lemma-stalk-of-pushforward}{{0AES}{source}{}{sourcepointer-4}{}}
\newlabel{lemma-stalk-of-pushforward@cref}{{[originalproposition][0][]0AES}{[1][0][]source}{}{}{}}
\makeatother
\crefname{originalproposition}{original proposition}{original propositions}
\Crefname{originalproposition}{Original proposition}{Original propositions}
\makeatletter
\newlabel{more-algebra-proposition-equivalence}{{05ER}{source}{}{sourcepointer-5}{}}
\newlabel{more-algebra-proposition-equivalence@cref}{{[originalproposition][0][]05ER}{[1][0][]source}{}{}{}}
\makeatother
\crefname{originalproposition}{original proposition}{original propositions}
\Crefname{originalproposition}{Original proposition}{Original propositions}
\makeatletter
\newlabel{more-algebra-remark-glueing-data}{{05EL}{source}{}{sourcepointer-6}{}}
\newlabel{more-algebra-remark-glueing-data@cref}{{[originalproposition][0][]05EL}{[1][0][]source}{}{}{}}
\makeatother
\crefname{originalproposition}{original proposition}{original propositions}
\Crefname{originalproposition}{Original proposition}{Original propositions}
\makeatletter
\newlabel{schemes-section-quasi-coherent-affine}{{01I6}{source}{}{sourcepointer-7}{}}
\newlabel{schemes-section-quasi-coherent-affine@cref}{{[originalproposition][0][]01I6}{[1][0][]source}{}{}{}}
\makeatother
\crefname{originalproposition}{original proposition}{original propositions}
\Crefname{originalproposition}{Original proposition}{Original propositions}
\makeatletter
\newlabel{situation-formal-glueing}{{0AEV}{source}{}{sourcepointer-8}{}}
\newlabel{situation-formal-glueing@cref}{{[originalproposition][0][]0AEV}{[1][0][]source}{}{}{}}
\makeatother
\title{Formal glueing is an equivalence for quasi-coherent modules with quasi-inverse given by the reverse functor}
\author{The Stacks Project Authors}
\date{stacks-project:spaces-pushouts.tex:local-snapshot-20261012; GNU FDL 1.2 or later}
\begin{document}
\maketitle
\noindent\textit{Editorial note.} Statement span is lines 2122-2127 of sources/stacks-spaces-pushouts/pinned-original-source.tex; proof span is lines 2129-2241, the complete author proof (111 body lines).\par\medskip
\noindent\textit{Editorial note.} Navigation stubs: 8 cited-but-not-included labels are supplied as non-load-bearing pointers with their pinned tags where the snapshot tag index has them.\par\medskip
\noindent\textit{Editorial note.} extra\_packages is empty on purpose: the retained preamble already loads xy with options and a second declaration clashes in the assembler wrapper.\par\medskip
\noindent\textit{Editorial note.} Text is retained verbatim, including the author's own typos and placeholders.\par\medskip
\noindent\textit{Editorial note.} Queue index 4 of queue\_spaces-pushouts.json (pb=4152, stmt 2122-2127, proof 2129-2241).\par\medskip
\noindent\textit{Editorial note.} spaces-pushouts wave-1: fresh adjudication accept (2026-10-14); six-check dispatched pool (§2 row appended 02:4x).\par\medskip
\noindent\textit{Allowed standard prerequisites (editorial).}\begin{itemize}
\item Explicit prerequisite boundary. The cited inputs are prerequisites; the argument in the retained span is not replaced by any cited result.
\end{itemize}\medskip
\ifdefined\hypertarget\hypertarget{sourcepointer-1}{}\fi
\noindent\textit{Source reference.} Original proposition 0AEW: sources/stacks-spaces-pushouts/pinned-original-source.tex, source lines 2073--2075. Cited by the retained statement or proof but not part of the collected spans; supplied only as a navigation pointer. The argument does not depend on its content.\par\medskip
\ifdefined\hypertarget\hypertarget{sourcepointer-2}{}\fi
\noindent\textit{Source reference.} Original proposition 0AEX: sources/stacks-spaces-pushouts/pinned-original-source.tex, source lines 2082--2084. Cited by the retained statement or proof but not part of the collected spans; supplied only as a navigation pointer. The argument does not depend on its content.\par\medskip
\ifdefined\hypertarget\hypertarget{sourcepointer-3}{}\fi
\noindent\textit{Source reference.} Original proposition 0AER: sources/stacks-spaces-pushouts/pinned-original-source.tex, source lines 1882--1896. Cited by the retained statement or proof but not part of the collected spans; supplied only as a navigation pointer. The argument does not depend on its content.\par\medskip
\ifdefined\hypertarget\hypertarget{sourcepointer-4}{}\fi
\noindent\textit{Source reference.} Original proposition 0AES: sources/stacks-spaces-pushouts/pinned-original-source.tex, source lines 1931--1944. Cited by the retained statement or proof but not part of the collected spans; supplied only as a navigation pointer. The argument does not depend on its content.\par\medskip
\ifdefined\hypertarget\hypertarget{sourcepointer-5}{}\fi
\noindent\textit{Source reference.} Original proposition 05ER: sources/stacks-spaces-pushouts/pinned-original-source.tex, source lines 2241--2241. Cited by the retained statement or proof but not part of the collected spans; supplied only as a navigation pointer. The argument does not depend on its content.\par\medskip
\ifdefined\hypertarget\hypertarget{sourcepointer-6}{}\fi
\noindent\textit{Source reference.} Original proposition 05EL: sources/stacks-spaces-pushouts/pinned-original-source.tex, source lines 2241--2241. Cited by the retained statement or proof but not part of the collected spans; supplied only as a navigation pointer. The argument does not depend on its content.\par\medskip
\ifdefined\hypertarget\hypertarget{sourcepointer-7}{}\fi
\noindent\textit{Source reference.} Original proposition 01I6: sources/stacks-spaces-pushouts/pinned-original-source.tex, source lines 2241--2241. Cited by the retained statement or proof but not part of the collected spans; supplied only as a navigation pointer. The argument does not depend on its content.\par\medskip
\ifdefined\hypertarget\hypertarget{sourcepointer-8}{}\fi
\noindent\textit{Source reference.} Original proposition 0AEV: sources/stacks-spaces-pushouts/pinned-original-source.tex, source lines 2046--2060. Cited by the retained statement or proof but not part of the collected spans; supplied only as a navigation pointer. The argument does not depend on its content.\par\medskip
\begin{proposition}
\label{proposition-formal-glueing-modules}
In Situation \ref{situation-formal-glueing} the functor
(\ref{equation-formal-glueing-modules}) is an equivalence
with quasi-inverse given by (\ref{equation-reverse}).
\end{proposition}

\begin{proof}
We first treat the special case where $X$ and $Y$ are affine schemes
and where the morphism $f$ is flat. Say $X = \Spec(R)$ and $Y = \Spec(S)$.
Then $f$ corresponds to a flat ring map $R \to S$. Moreover, $Z \subset X$
is cut out by a finitely generated ideal $I \subset R$. Choose generators
$f_1, \ldots, f_t \in I$. By the description of quasi-coherent modules
in terms of modules
(Schemes, Section \ref{schemes-section-quasi-coherent-affine}),
we see that the category $\textit{QCoh}(Y \to X, Z)$
is canonically equivalent to the category
$\text{Glue}(R \to S, f_1, \ldots, f_t)$
of More on Algebra, Remark \ref{more-algebra-remark-glueing-data}
such that the functors
(\ref{equation-formal-glueing-modules}) and (\ref{equation-reverse})
correspond to the functors $\text{Can}$ and $H^0$.
Hence the result follows from
More on Algebra, Proposition \ref{more-algebra-proposition-equivalence}
in this case.

\medskip\noindent
We return to the general case.
Let $\mathcal{F}$ be a quasi-coherent module on $X$.
We will show that
$$
\alpha :
\mathcal{F}
\longrightarrow
\Ker\left(j_*\mathcal{F}|_U \oplus f_*f^*\mathcal{F} \to
(f \circ j')_*f^*\mathcal{F}|_V\right)
$$
is an isomorphism. Let $(\mathcal{H}, \mathcal{G}, \varphi)$
be an object of $\QCoh(Y \to X, Z)$. We will show that
$$
\beta :
f^*\Ker\left(
j_*\mathcal{H} \oplus f_*\mathcal{G} \to (f \circ j')_*\mathcal{G}|_V
\right)
\longrightarrow
\mathcal{G}
$$
and
$$
\gamma :
j^*\Ker\left(
j_*\mathcal{H} \oplus f_*\mathcal{G} \to (f \circ j')_*\mathcal{G}|_V
\right)
\longrightarrow
\mathcal{H}
$$
are isomorphisms. To see these statements are true it suffices to
look at stalks. Let $\overline{y}$ be a geometric point of $Y$ mapping
to the geometric point $\overline{x}$ of $X$.

\medskip\noindent
Fix an object $(\mathcal{H}, \mathcal{G}, \varphi)$ of $\QCoh(Y \to X, Z)$.
By Lemma \ref{lemma-stalk-formal-glueing}
and a diagram chase (omitted) the canonical map
$$
\Ker(j_*\mathcal{H} \oplus f_*\mathcal{G} \to
(f \circ j')_*\mathcal{G}|_V)_{\overline{x}}
\longrightarrow
\Ker(
j_*\mathcal{H}_{\overline{x}} \oplus \mathcal{G}_{\overline{y}}
\to
j'_*\mathcal{G}_{\overline{y}}
)
$$
is an isomorphism.

\medskip\noindent
In particular, if $\overline{y}$ is a geometric point of $V$, then
we see that $j'_*\mathcal{G}_{\overline{y}} = \mathcal{G}_{\overline{y}}$
and hence that this kernel is equal to $\mathcal{H}_{\overline{x}}$.
This easily implies that $\alpha_{\overline{x}}$, $\beta_{\overline{x}}$,
and $\beta_{\overline{y}}$ are isomorphisms in this case.

\medskip\noindent
Next, assume that $\overline{y}$ is a point of $f^{-1}Z$.
Let $I_{\overline{x}} \subset \mathcal{O}_{X, \overline{x}}$,
resp.\ $I_{\overline{y}} \subset \mathcal{O}_{Y, \overline{y}}$
be the stalk of the ideal cutting out $Z$, resp.\ $f^{-1}Z$.
Then $I_{\overline{x}}$ is a finitely generated ideal,
$I_{\overline{y}} = I_{\overline{x}}\mathcal{O}_{Y, \overline{y}}$,
and $\mathcal{O}_{X, \overline{x}} \to \mathcal{O}_{Y, \overline{y}}$
is a flat local homomorphism inducing an isomorphism
$\mathcal{O}_{X, \overline{x}}/I_{\overline{x}} =
\mathcal{O}_{Y, \overline{y}}/I_{\overline{y}}$.
At this point we can bootstrap using the diagram of categories
$$
\xymatrix{
\QCoh(\mathcal{O}_X) \ar[r]_-{(\ref{equation-formal-glueing-modules})} \ar[d] &
\QCoh(Y \to X, Z) \ar[d] \ar@/_2pc/[l]^{(\ref{equation-reverse})} \\
\text{Mod}_{\mathcal{O}_{X, \overline{x}}} \ar[r]^-{\text{Can}} &
\text{Glue}(\mathcal{O}_{X, \overline{x}} \to \mathcal{O}_{Y, \overline{y}},
f_1, \ldots, f_t) \ar@/^2pc/[l]_{H^0}
}
$$
Namely, as in the first paragraph of the proof we identify
$$
\text{Glue}(\mathcal{O}_{X, \overline{x}} \to \mathcal{O}_{Y, \overline{y}},
f_1, \ldots, f_t)
=
\QCoh(\Spec(\mathcal{O}_{Y, \overline{y}}) \to
\Spec(\mathcal{O}_{X, \overline{x}}), V(I_{\overline{x}}))
$$
The right vertical functor is given by pullback, and it is clear that
the inner square is commutative. Our computation of the stalk of the
kernel in the third paragraph of the proof combined with
Lemma \ref{lemma-stalk-of-pushforward} implies that
the outer square (using the curved arrows) commutes. Thus we
conclude using the case of a flat morphism of affine schemes
which we handled in the first paragraph of the proof.
\end{proof}

\end{document}
